\documentclass{article}

\usepackage{corl_2026}
\usepackage{amsmath}
\usepackage{amssymb}
\usepackage{booktabs}
\usepackage{graphicx}
\usepackage{xcolor}
\usepackage{multirow}

% Placeholder command — fill in before submission
\newcommand{\todo}[1]{\textcolor{red}{[TODO: #1]}}

\title{Active Morphological Probing: Terrain-Aware Locomotion Without Depth Perception}

\author{
  % NOTE: anonymized for initial submission
}

\begin{document}
\maketitle

%===============================================================================

\begin{abstract}
Wheeled-legged robots gives the versatility of legged locomotion on unstructured terrain and the
speed and efficiency of wheeled movement on flat ground. The existing robot design actively uses all 
degrees of freedom to move across different terrains
incorporates 
by adding wheel to the feet of the robots, which adds four extra degrees of freedom at the feet 
of the quadruped robots that makes it less power efficient. In this paper, we present an open source
morphologically adaptive robot that actively reconfigures between wheeled and legged locomotion to move
effortlessly between varied terrains. 

but typically require depth cameras and complex algorithms
to determine when to reconfigure. We argue that robot design choices directly determine
algorithmic complexity a robot built with learnability in mind can reduce a hard
perception-and-control problem to a simple learned policy. We present AWM (All-Wheel-Morph),
a low-cost open-source wheeled-legged robot where each wheel extends into a leg, designed
to make morphological adaptation learnable: actuators with built-in torque observability,
a continuous morphology range, and sim-to-real transferable dynamics. Leveraging torque
feedback natively available from AWM's actuators, we propose \textbf{Active Morphological
Probing} a GRU policy that physically senses terrain by extending legs and reading
contact torque, adapting morphology without any depth camera. Without exteroception, our
policy matches camera-equipped baselines in terrain traversal while a fixed-morphology
wheels-only policy plateaus at half the curriculum difficulty regardless of training
duration, confirming that adaptive morphology is the key enabler. Our approach requires
no teacher-student distillation, no privileged training information, and deploys as a
single network. We fully open-source AWM including hardware design, bill of materials,
simulation assets, training code, and trained policies so that anyone can build, train,
and deploy a morphologically adaptive robot with off-the-shelf components under \$2,500.
\end{abstract}

\keywords{morphological adaptation, legged locomotion, proprioception, torque sensing, sim-to-real}

%===============================================================================

\section{Introduction}

Traversing unstructured terrain stairs, curbs, rough ground remains an open problem
for wheeled robots. The dominant solution couples a depth camera with a learned policy that
reads a height map of the upcoming surface and decides how to reconfigure or accelerate
\citep{lee2020learning, miki2022learning}. This works, but it couples the robot's
intelligence to a sensor that adds cost, weight, and failure modes: cameras saturate in
bright light, fail on reflective or textureless surfaces, and require careful extrinsic
calibration that drifts over time.

We take a different starting point. Rather than asking what sensor can tell the robot
about terrain, we ask what physical design makes terrain information natively
available. Our answer is morphological probing: a robot that extends its legs into
contact with the surface and reads the resulting joint torques. Stairs feel different from
flat ground through torque alone --- a stair edge loads the leg abruptly, producing a
distinctive current spike that a GRU can learn to recognize over time. No depth image
required.

This idea only works if the robot is designed for it. A leg with a high-ratio gearbox in
position-control mode cannot sense torque: gear friction holds most of the load, and motor
current does not reflect contact force. We present AWM (All-Wheel-Morph), a wheeled-legged
platform specifically designed to make torque sensing reliable: current-based position
control (operating mode 5) on all leg actuators, calibrated to within 2\% across 18 trials,
producing torque estimates that map directly to the simulation's observation space without
additional processing.

We train a GRU policy over a 25-dimensional proprioceptive observation that includes
four normalized leg torques but no exteroception. The result is \textbf{Active Morphological
Probing (AMP)}: the policy learns to extend legs as terrain probes, read the torque
response, and adapt morphology accordingly. In simulation on a staircase curriculum,
AMP matches the terrain traversal performance of a camera-equipped baseline (GRU +
terrain\_scan: $0.72$ vs.\ AMP: $0.72 \pm 0.03$), while a wheels-only policy with
identical training plateaus at $0.42$, confirming that adaptive morphology --- not just
a better policy --- is the key enabler.

Our contributions are:
\begin{enumerate}
    \item \textbf{AWM}: an open-source wheeled-legged robot under \$2{,}500 with calibrated
    torque sensing via current-based position control, designed for sim-to-real transfer.
    \item \textbf{Active Morphological Probing}: a GRU policy that physically senses
    terrain through leg contact torques, requiring no depth camera, teacher network, or
    privileged training signal.
    \item Simulation experiments showing AMP matches camera-equipped baselines and
    dramatically outperforms fixed-morphology wheels-only locomotion. \todo{add: real-robot
    sim-to-real validation}
    \item Full open-source release: CAD, bill of materials, simulation assets, training
    code, and trained weights.
\end{enumerate}

%===============================================================================

\section{Related Work}

\paragraph{Terrain-aware locomotion with exteroception.}
\citet{miki2022learning} show that an attention-based recurrent encoder fusing proprioception
with a depth-camera height map enables robust locomotion in unstructured outdoor environments.
\citet{agarwal2022legged} achieve stair and gap traversal using only a front-facing depth
camera in an end-to-end policy. \citet{lee2024wheeled} extend perceptive locomotion to
wheeled-legged ANYmal, validating km-scale navigation in urban environments with privileged
height-map training. All of these require exteroceptive sensors at deployment; AMP matches
their terrain performance without any depth sensor by using leg torques instead.

\paragraph{Proprioceptive and blind locomotion.}
\citet{lee2020learning} showed that an LSTM over proprioceptive history enables
challenging-terrain quadruped locomotion without any camera.
RMA~\citep{kumar2021rma} extends this with a two-phase approach: a base policy trained
with privileged environment information and an adaptation module that infers those factors
from proprioceptive history at deployment.
DreamWaQ~\citep{nahrendra2023dreamwaq} learns an implicit terrain representation through
a latent state estimator without exteroception.
\citet{ji2024hybrid} use a GRU-based internal model with joint encoders and IMU only,
estimating terrain state from the robot's own dynamic response.
\citet{margolis2024rapid} achieve 3.9\,m/s blind locomotion on MIT Mini Cheetah with
a gait-conditioned proprioceptive policy.
All of these reactive approaches detect terrain \emph{after} contact; AMP deliberately
extends legs \emph{before} committing to traversal, gathering terrain information actively.

\paragraph{Morphologically adaptive and wheeled-legged robots.}
Wheeled-legged platforms combine rolling efficiency with stepping capability for versatile
terrain coverage \citep{bjelonic2019keep, bjelonic2021whole, klemm2019ascento}.
\citet{rudin2024blind} demonstrate blind stair climbing on the wheeled-legged Ascento
using asymmetric actor-critic RL, achieving sim-to-real transfer without exteroception.
These works treat morphology as fixed or switch between discrete modes via trajectory
optimization. AWM learns a \emph{continuous} morphological adaptation policy driven
entirely by sensory feedback, without any explicit mode-switching logic.

\paragraph{Curriculum learning for locomotion RL.}
\citet{rudin2022learning} established the massively parallel terrain curriculum paradigm
that Isaac Lab descends from, training locomotion policies in minutes on thousands of
parallel environments with automatic difficulty progression. Walk These Ways~\citep{margolis2022walk}
extends this with gait command conditioning for behavioral generalization across locomotion
styles. AWM's terrain curriculum follows this design directly.

\paragraph{Joint torques and contact sensing.}
Joint torque sensing is used for contact detection, impedance control, and collision
identification in manipulation \citep{haddadin2017robot}, and for foot contact detection
in legged locomotion. \citet{higuera2022traversability} use foot-mounted force probes for
terrain traversability assessment in a planning module separate from the locomotion policy.
To our knowledge, no prior work feeds raw joint torques as a direct observation into
a locomotion RL policy. We show that a GRU over 4-dimensional leg torque history can
replace a 35-dimensional terrain height map for morphological adaptation decisions.

%===============================================================================

\section{The AWM Platform}
\label{sec:platform}

\begin{figure}[t]
    \centering
    \todo{Figure: robot photo (real) + CAD render side by side, with leg extended and
    retracted positions labeled}
    \caption{AWM (All-Wheel-Morph). Each of the four wheels integrates a leg actuator
    whose rotation extends the wheel hub into ground contact. At left: legs retracted
    (wheel mode). At right: legs extended (probing/stair mode). All hardware is
    open-sourced under \$2{,}500.}
    \label{fig:robot}
\end{figure}

\subsection{Mechanical Design}

AWM is a four-wheeled robot in which each wheel hub is coupled to a rotary leg actuator.
Rotating the actuator extends the wheel outward and downward, transitioning from pure
rolling to a hybrid wheel-leg contact. The morphology range is continuous: from fully
retracted (wheels contact ground, legs clear) to fully extended (legs contact ground,
providing torque measurement). This design choice --- continuous rather than discrete
morphology --- is deliberate: it gives the policy a smooth action space and allows
partial extension as a probing strategy without committing to a full stance change.

Wheels are driven by XM430-W210-R actuators (RS-485 bus, velocity control). Legs are
driven by XM430-W350-T actuators (TTL bus, current-based position control, stall torque
4.1\,Nm at 12\,V). Four slip rings (6-wire, 10\,A rated) provide power and signal
continuity through leg rotation. The robot is designed for straightforward assembly from
off-the-shelf components; the full bill of materials is provided in the supplementary
materials.

\subsection{Torque Sensing via Actuation Mode}
\label{sec:torque_sensing}

The key hardware insight is that torque sensing quality depends critically on actuator
operating mode. In standard position control (mode 3), the high gear ratio (353:1) of the
XM430-W350-T causes gear friction to hold most external load --- Present Current reflects
the motor's effort to maintain position against friction, not against terrain. We switch
all leg actuators to \textbf{current-based position control (mode 5)}: the motor
actively maintains a goal position by pushing against terrain, and Present Current directly
reflects terrain reaction force.

We calibrate the torque constant empirically: a rigid test arm at three distances (40,
80, 120\,mm) with two known loads (500\,g, 1\,kg) across 18 trials yields
$k_\tau = 3.83$\,Nm/A (std $< 0.12$\,Nm/A across all conditions). The conversion is:
\begin{equation}
    \tau_\text{Nm} = I_\text{LSB} \times 2.69 \times 10^{-3} \times 3.83
    \label{eq:torque}
\end{equation}
where $I_\text{LSB}$ is the raw Present Current register value and $2.69$\,mA/LSB is
the unit from the actuator datasheet. This pipeline runs at 60\,Hz on the Jetson Orin
Nano Super, matching the simulation policy rate.

\subsection{Open-Source Release}

AWM is fully open-sourced: CAD files, bill of materials, simulation assets (Isaac Lab USD),
training code, and trained policy weights are released at \todo{URL}. The complete
platform can be assembled from off-the-shelf components for under \$2{,}500, enabling
any lab to replicate our experiments.

%===============================================================================

\section{Active Morphological Probing}
\label{sec:method}

\subsection{Problem Formulation}

We formulate locomotion as a Markov Decision Process (MDP) with state $s_t$, action
$a_t$, and reward $r_t$. The robot follows a commanded velocity $(v_x, \omega_z)$. The
policy $\pi_\theta$ maps an observation history to actions: wheel velocities for the four
drive motors and leg position targets for the four leg actuators.

\subsection{Observation Space}

The AMP policy receives a 25-dimensional proprioceptive observation with no exteroceptive
signal:

\begin{table}[h]
\centering
\small
\begin{tabular}{lcc}
\toprule
Observation & Dims & Source \\
\midrule
Commanded velocity $(v_x, \omega_z)$ & 2 & Joystick / planner \\
Base linear velocity $v_x$ & 1 & Wheel odometry \\
Base angular velocity $\omega_z$ & 1 & IMU (gyro $z$) \\
Wheel velocities & 4 & Drive motor encoders \\
Leg positions & 4 & Leg motor encoders \\
Previous leg actions & 4 & Policy memory \\
Projected gravity & 3 & IMU (accelerometer) \\
Progress--slip history & 2 & Computed \\
\textbf{Leg torques} & \textbf{4} & \textbf{Mode-5 current (Eq.~\ref{eq:torque})} \\
\midrule
\textbf{Total} & \textbf{25} & \\
\bottomrule
\end{tabular}
\caption{AMP observation space. Leg torques (bold) replace the 35-dimensional
terrain\_scan from the camera baseline, reducing observation size by 31 dimensions while
enabling terrain sensing through physical contact.}
\label{tab:obs}
\end{table}

\paragraph{Torque sensing via PD loading.}
The leg torque signal arises from a deliberate interaction between the policy and the
PD controller. At each timestep, the policy outputs a leg position target $q^\text{des}$,
and the actuator applies torque according to:
\begin{equation}
    \tau = k_p (q^\text{des} - q) + k_d (0 - \dot{q})
    \label{eq:pd}
\end{equation}
On flat ground the leg reaches $q^\text{des}$ easily and $\tau$ remains small. When
the leg contacts a stair edge, terrain physically prevents $q$ from reaching $q^\text{des}$
--- the position error $(q^\text{des} - q)$ stays large, the motor draws proportionally
more current, and $\tau$ spikes. This spike enters $x_{t+1}$ as the leg torque observation.
Critically, the policy \emph{causes} this signal by commanding a target the terrain cannot
satisfy --- probing terrain stiffness rather than passively measuring contact. The GRU
integrates this torque history over $\approx 0.8$\,s to distinguish terrain types: a stair
produces a sustained asymmetric spike pattern across the four legs, while flat ground
produces near-zero symmetric readings.

Leg torques are normalized by the actuator effort limit ($2.70$\,Nm) and clamped to
$[-2, 2]$ to handle transient impact spikes. Gaussian noise (std $= 0.05$) is added
during training for sim-to-real robustness.

The camera baseline (GRU + terrain\_scan) uses a 56-dimensional observation including a
$7\times5$ height grid from a ray-caster sensor (in simulation) or ZED Mini depth camera
(on the real robot), with matched Gaussian noise (std $= 0.02$, 2\,cm depth noise model).

\subsection{Policy Architecture}

Both AMP and the camera baseline use the same GRU architecture: a single-layer GRU
with 256 hidden units followed by a two-layer MLP actor head (hidden dims [256, 128], ELU
activations). The GRU processes the full observation at each timestep and maintains hidden
state across the episode, allowing the policy to integrate torque history for terrain
identification. The recurrent rollout length is 48 steps ($\approx 0.8$\,s at 60\,Hz).
A separate critic network shares the same structure. Action noise is initialized at
$\sigma = 0.3$ (scalar, not state-dependent).

\subsection{Training}
\label{sec:training}

We train using Proximal Policy Optimization (PPO)~\citep{schulman2017proximal} in
Isaac Lab~\citep{mittal2023orbit} with 1024 parallel environments. The terrain curriculum
consists of flat ground (50\%) and random uniform staircases (50\%, step height
0.05--0.20\,m, step width 0.25--0.35\,m). The curriculum difficulty (terrain\_levels,
$\in [0, 1]$) advances when the mean episode reward exceeds a threshold, and the metric
at convergence is our primary performance measure.

The reward function encourages velocity tracking and penalizes tilt, wheel slip,
and excessive joint velocities:
\begin{equation}
r_t = r_\text{vel} + r_\text{yaw} + r_\text{leg} - c_\text{tilt} - c_\text{slip}
      - c_\text{joint} - c_\text{action}
\end{equation}
A morphology-specific bonus $r_\text{leg}$ rewards leg extension on rough terrain and
penalizes retracted legs when stuck, encouraging the emergent probing behavior. All
policies are trained for 30{,}000 iterations ($\approx 1.47 \times 10^9$ environment steps).

%===============================================================================

\section{Experiments}
\label{sec:experiments}

We evaluate three questions: (1) Does torque sensing match camera performance in
simulation? (2) Is the morphology specifically contributing, or just having legs?
(3) Does the policy transfer to the real robot?

\subsection{Simulation Comparison}
\label{sec:sim_comparison}

We compare four conditions trained on the same curriculum for the same number of
iterations:

\begin{table}[h]
\centering
\begin{tabular}{lcccc}
\toprule
Method & Architecture & Obs dims & terrain\_levels & Seeds \\
\midrule
WheelsOnly           & MLP  & 21  & $0.42$           & 1 \\
MLP + terrain\_scan  & MLP  & 56  & $0.68 \pm 0.02$  & 3 \\
GRU + terrain\_scan & GRU & 56  & $0.72$           & 1 \\
\textbf{AMP (ours)}  & GRU & \textbf{25}  & $\mathbf{0.72 \pm 0.03}$ & \textbf{4} \\
\bottomrule
\end{tabular}
\caption{Terrain curriculum level at convergence (mean $\pm$ std across seeds, last 10
evaluations per run). AMP matches the GRU camera baseline while using 31 fewer
observation dimensions and no depth sensor.}
\label{tab:sim_results}
\end{table}

Three findings stand out. First, AMP ($0.72 \pm 0.03$) matches GRU + terrain\_scan
($0.72$) on the apples-to-apples comparison: same architecture, same training, only
the sensing modality differs. Second, morphology is essential: WheelsOnly plateaus at
$0.42$ regardless of training duration, a gap of $0.30$ terrain levels that no policy
improvement can close. Third, the recurrent architecture matters: MLP + terrain\_scan
($0.68$) underperforms GRU + terrain\_scan ($0.72$), confirming that temporal
integration of terrain signals --- whether from depth or torque --- is important.

\subsection{Ablation: What Does Torque Sensing Contribute?}
\label{sec:ablation}

\todo{Run: GRU + legs extended fixed (no torque sensing) condition. This isolates
whether torque specifically is doing the work vs.\ just having legs available. Expected
result: fixed-leg condition underperforms AMP, showing torque information --- not just
morphology presence --- drives the gain over WheelsOnly.}

\begin{table}[h]
\centering
\begin{tabular}{lcc}
\toprule
Method & Legs & terrain\_levels \\
\midrule
WheelsOnly                     & Retracted (fixed) & $0.42$ \\
GRU + fixed legs (no torque)   & Extended (fixed)  & \todo{X.XX} \\
\textbf{AMP (ours)}            & Adaptive + torque & $0.72 \pm 0.03$ \\
\bottomrule
\end{tabular}
\caption{Ablation isolating the contribution of torque sensing. Fixed-legs condition
has legs extended but not torque-sensing: it cannot adapt morphology and receives no
torque input.}
\label{tab:ablation}
\end{table}

\subsection{Torque Signatures as Terrain Encoding}
\label{sec:torque_signatures}

\todo{Figure: real robot leg torque traces over time on (a) flat ground, (b) stair
approach and climb. Show that stairs produce a characteristic torque spike pattern that
is temporally structured and repeatable across trials. This is the physical justification
for why a GRU over torque history can substitute for a depth image.}

To understand \emph{why} AMP works, Figure~\todo{X} shows leg torque traces from the
real robot traversing flat ground and a three-step staircase. \todo{Write 2--3 sentences
describing the pattern after figure is available.}

\subsection{Sim-to-Real Transfer}
\label{sec:sim2real}

\todo{Run both AMP and GRU + terrain\_scan on the real robot. Report terrain\_levels
using the same metric as simulation (stair step height curriculum). Target: within 10\%
of sim performance for both conditions. Use ZED Mini IMU for both (same hardware),
ZED Mini depth for terrain\_scan condition only.}

\begin{table}[h]
\centering
\begin{tabular}{lcc}
\toprule
Method & Sim & Real \\
\midrule
GRU + terrain\_scan & $0.72$ & \todo{X.XX} \\
\textbf{AMP (ours)}  & $0.72 \pm 0.03$ & \todo{X.XX} \\
\bottomrule
\end{tabular}
\caption{Sim-to-real transfer. Both policies use identical hardware (ZED Mini for IMU);
terrain\_scan additionally uses ZED Mini depth. The clean hardware match isolates the
sensing modality as the only difference between conditions.}
\label{tab:sim2real}
\end{table}

\subsection{Failure Analysis}
\label{sec:failure}

\todo{After real-robot experiments: document conditions where AMP fails. Likely candidates:
(1) terrain step height exceeding leg range, (2) slippery surfaces where torque ambiguously
reflects slip vs.\ terrain resistance, (3) rapid step sequences that exceed GRU temporal
integration window.}

%===============================================================================

\section{Conclusion}

We presented Active Morphological Probing, an approach that uses a robot's own leg
actuators as terrain sensors. By switching leg actuators to current-based position control
and feeding normalized torque signals into a GRU policy, we match the terrain traversal
performance of a depth-camera-equipped baseline without any exteroceptive sensor. The
key insight is that robot design determines sensing capability: the same actuators that
create morphological adaptability also provide terrain information, at no additional
hardware cost.

The performance gap between WheelsOnly ($0.42$) and AMP ($0.72$) demonstrates that
adaptive morphology is the primary enabler of stair climbing performance --- a gap that
training time alone cannot close. We release AWM, the full hardware and software stack,
to make this a reproducible community benchmark for morphologically adaptive locomotion.

%===============================================================================

\clearpage
\acknowledgments{\todo{Add acknowledgments for camera-ready version.}}

\bibliography{references}

\end{document}
